% !TEX program = xelatex
% 第一讲：second_order.tex
% 二阶线性常微分方程：奇点、Frobenius 级数、第二个解
\documentclass[fontset=windows, chinese, 10pt, xcolor=dvipsnames]{beamer}
% Shared preamble for the chapter 07 standalone lectures.
% Each lecture sets \title and \subtitle after this file is read.

\usepackage{amsmath,amssymb,mathtools}
\usepackage{xeCJK}
\usepackage{zhnumber}
\usepackage{tikz}
\usepackage{tikz-3dplot}
\usetikzlibrary{arrows.meta,calc,positioning}

\usetheme{metropolis}
\usepackage{beamercolorthemeHistNanjing}

\newcommand{\red}[1]{\textcolor{red}{#1}}
\newcommand{\blue}[1]{\textcolor{blue}{#1}}
\newcommand{\green}[1]{\textcolor{green}{#1}}
\newcommand{\gray}[1]{\textcolor{gray}{#1}}
\newcommand{\orange}[1]{\textcolor{orange}{#1}}

\newcommand{\sssub}[1]{ {\mathrm{#1} }}  % script size subscript\\\
\newcommand{\funcd}[2]{ \frac{d{#1}}{d{#2}}} %function full derivative
\newcommand{\funcpd}[2]{ \frac{\partial{#1}}{\partial{#2}}} %function partial derivative
\newcommand{\fnald}[2] { \frac{\delta {#1}}{\delta {#2}}} % functional derivative

\renewcommand{\Re}{\operatorname{Re}}
\renewcommand{\Im}{\operatorname{Im}}
\def\ext{\sssub{ext}}
\def\x{\sssub{x}}
\def\xc{\sssub{xc}}
\def\H{\sssub{H}}
\def\Hx{\sssub{Hx}}
\def\Hxc{\sssub{Hxc}}
\def\c{\sssub{c}}
\def\s{\sssub{s}}

% \def\ext{\sssub{ext}}
% \def\ext{\sssub{ext}}
% \def\ext{\sssub{ext}}
% \def\ext{\sssub{ext}}
% \defsssub{ext}


% definitions of bold faced symbols
\def\brpppp{{\mathbf{r}^{\prime\prime\prime\prime}}}
\def\brppp{{\mathbf{r}^{\prime\prime\prime}}}
\def\brpp{{\mathbf{r}^{\prime\prime}}}
\def\brp{{\mathbf{r}^{\prime}}}
\def\bzp{{\mathbf{z}^{\prime}}}
\def\bxp{{\mathbf{x}^{\prime}}}
\def\tp{{{t}^{\prime}}}
\def\tpp{{{t}^{\prime\prime}}}
\def\tppp{{{t}^{\prime\prime\prime}}}

\def\tbr{{\tilde{\mathbf{r}}}}
\def\bk{{\mathbf{k}}}
\def\br{{\mathbf{r}}}
\def\bp{{\mathbf{p}}}
\def\bq{{\mathbf{q}}}
\def\bv{{\mathbf{v}}}
\def\bz{{\mathbf{z}}}
\def\mbf{{\mathbf{f}}}

\def\bx{{\mathbf{x}}}
\def\bR{{\mathbf{R}}}
\def\bM{{\mathbf{M}}}
\def\bP{{\mathbf{P}}}
\def\bT{{\mathbf{T}}}
\def\bK{{\mathbf{K}}}
\def\bA{{\mathbf{A}}}
\def\bB{{\mathbf{B}}}
\def\bD{{\mathbf{D}}}
\def\bE{{\mathbf{E}}}
\def\bF{{\mathbf{F}}}
\def\bG{{\mathbf{G}}}
\def\bH{{\mathbf{H}}}
\def\bI{{\mathbf{I}}}
\def\bJ{{\mathbf{J}}}
\def\bP{{\mathbf{P}}}
\def\bS{{\mathbf{S}}}
\def\bV{{\mathbf{V}}}
\def\bX{{\mathbf{X}}}
\def\bY{{\mathbf{Y}}}
\def\dd{{\mathop{}\!\mathrm{d}}}
\def\dr{{\!\!\dd\br\,}}
\def\dx{{\!\!\dd\bx\,}}
\def\dy{{\!\!\dd\by\,}}
\def\dk{{\!\!\dd\bk\,}}
\def\dG{{\!\!\dd\bG\,}}


\def\half{{\frac{1}{2}}}
\def\tr{{\mathrm{tr}}}


\def\Arg{{\mathrm{Arg}\;}}
\def\arg{{\mathrm{arg}\;}}
\def\Res{{\mathrm{Res}\;}}

% unit-vector basis (bold under XeLaTeX; \mathbf can look non-bold)
\newcommand{\ue}{\boldsymbol{e}}
\newcommand{\un}{\hat{\boldsymbol{n}}}


\author{罗凯}
\institute{南京理工大学物理学院}
\date{\zhdigits{\the\year},\zhganzhinian{\the\year}年秋}

\usebackgroundtemplate{%
\begin{tikzpicture}[remember picture,overlay]
  \node[opacity=0.05, at=(current page.center)] {
    \includegraphics[width=0.6\paperwidth]{njust-logo.jpeg}
  };
\end{tikzpicture}}


\title{二阶微分方程}
\subtitle{第一讲：奇点与级数解}

\begin{document}

\begin{frame}
  \titlepage
\end{frame}

\begin{frame}{本讲提纲}
  分离变量把偏微分方程收成常微分方程。二阶线性方程通常没有初等通解，先看奇点，再在常点或正则奇点附近作幂级数。
  \begin{enumerate}
    \item 常点、正则奇点、非正则奇点；
    \item 常点附近的 Frobenius 级数，用谐振子认出正弦和余弦；
    \item 正则奇点：贝塞尔方程的第一个解；
    \item 朗斯基行列式给出第二个线性无关解。
  \end{enumerate}
\end{frame}

% =============================================================
\section{奇点}

\begin{frame}{标准形式}
  二阶线性齐次方程写成
  \begin{equation*}
    y''+p(x)\,y'+q(x)\,y=0.
  \end{equation*}
  右端换成已知函数 $g(x)$，就是非齐次方程。叠加只对齐次方程成立：两个解的线性组合仍是解。
  \vspace{0.4em}
  通解含两个任意常数。找到两个线性无关的解，通解就是它们的组合。
\end{frame}

\begin{frame}{常点}
  若 $p(x)$ 和 $q(x)$ 在 $x_0$ 处都有限，称 $x_0$ 为\textbf{常点}（ordinary point）。
  \vspace{0.45em}
  常点附近两个解都解析，可以写成普通幂级数
  \begin{equation*}
    y(x)=\sum_{k=0}^{\infty}a_k(x-x_0)^k.
  \end{equation*}
  收敛半径至少延伸到最近的奇点。
\end{frame}

\begin{frame}{奇点}
  $p$ 或 $q$ 在 $x_0$ 处发散，$x_0$ 就是\textbf{奇点}（singular point）。再看发散有多快。
  \vspace{0.35em}
  若 $(x-x_0)p(x)$ 和 $(x-x_0)^2 q(x)$ 在 $x_0$ 处仍有限，称 $x_0$ 为\textbf{正则奇点}（regular singular point）。否则是\textbf{非正则奇点}。
  \vspace{0.35em}
  正则奇点附近，至少有一个解形如 $(x-x_0)^s$ 乘上幂级数，$s$ 不必是整数。
\end{frame}

\begin{frame}{例：勒让德方程}
  \[
    (1-x^2)y''-2xy'+\ell(\ell+1)y=0.
  \]
  除以 $1-x^2$，
  \[
    p(x)=-\frac{2x}{1-x^2},\qquad
    q(x)=\frac{\ell(\ell+1)}{1-x^2}.
  \]
  在 $x=\pm 1$ 发散。在 $x=1$，
  \[
    (x-1)p(x)=\frac{2x}{1+x},\qquad
    (x-1)^2 q(x)=-(x-1)\frac{\ell(\ell+1)}{1+x},
  \]
  都有限。$x=-1$ 同样。所以 $\pm 1$ 是正则奇点，其余有限点都是常点。
\end{frame}

\begin{frame}{例：贝塞尔方程}
  \[
    x^2 y''+xy'+(x^2-n^2)y=0.
  \]
  除以 $x^2$，
  \[
    p(x)=\frac{1}{x},\qquad
    q(x)=1-\frac{n^2}{x^2}.
  \]
  $xp=1$，$x^2 q=x^2-n^2$，在 $x=0$ 有限。原点是正则奇点。
  \vspace{0.35em}
  $x\neq 0$ 的有限点上 $p$、$q$ 都有限，都是常点。
\end{frame}

\begin{frame}{几种方程的奇点}
  \begin{center}
  \renewcommand{\arraystretch}{1.25}
  {\small
  \begin{tabular}{lc}
    方程 & 有限处的正则奇点 \\ \hline
    超几何 & $0,\,1$ \\
    勒让德、切比雪夫 & $-1,\,1$ \\
    合流超几何、拉盖尔、贝塞尔 & $0$
  \end{tabular}}
  \end{center}
  \vspace{0.35em}
  无穷远要另看。令 $x=1/z$，把方程改写成 $w(z)=y(1/z)$ 的方程，再对 $z=0$ 用同一套判据。
  超几何、勒让德、切比雪夫在无穷远是正则奇点；后三种在无穷远是非正则奇点。
\end{frame}

\begin{frame}{两种可以直接积分的方程}
  $y''=f(x)$ 积分两次：
  \[
    y=\int^{x}\!\!\int^{t}f(s)\,\mathrm{d}s\,\mathrm{d}t+C_1 x+C_2.
  \]
  $y''=f(x,y')$ 不含 $y$。令 $u=y'$，得到一阶方程 $u'=f(x,u)$。解出 $u=\varphi(x,C_1)$ 后再积分，
  \[
    y=\int\varphi(x,C_1)\,\mathrm{d}x+C_2.
  \]
  系数依赖 $y$ 本身时，这样降不下去，下面改用级数。
\end{frame}

% =============================================================
\section{常点：谐振子}

\begin{frame}{Frobenius 级数}
  在 $x_0=0$ 附近设
  \begin{equation*}
    y(x)=\sum_{j=0}^{\infty}a_j x^{s+j},\qquad a_0\neq 0.
  \end{equation*}
  $s$ 和 $a_j$ 都待定。求导，
  \begin{align*}
    y'&=\sum_{j=0}^{\infty}a_j(s+j)x^{s+j-1},\\
    y''&=\sum_{j=0}^{\infty}a_j(s+j)(s+j-1)x^{s+j-2}.
  \end{align*}
  代进微分方程，每一幂次的系数必须为零。最低次幂给出 $s$ 的\textbf{判别方程}（indicial equation）。
\end{frame}

\begin{frame}{谐振子}
  \[
    y''+\omega^2 y=0.
  \]
  $p=0$、$q=\omega^2$，原点是常点，但我们仍保留 $s$，用来对照后面的正则奇点。代入后最低次是 $x^{s-2}$：
  \[
    a_0\,s(s-1)=0.
  \]
  $a_0\neq 0$，所以 $s=0$ 或 $s=1$。解从 $x^0$ 或 $x^1$ 开始，和解析性一致。
\end{frame}

\begin{frame}{下一项与递推}
  $x^{s-1}$ 的系数是 $a_1(s+1)s=0$。
  $s=1$ 时必须 $a_1=0$；$s=0$ 时 $a_1$ 不受限制。
  \vspace{0.3em}
  更高次给出递推
  \begin{equation*}
    a_{j+2}=-\frac{\omega^2}{(s+j+2)(s+j+1)}\,a_j,
    \qquad j\geq 0.
  \end{equation*}
  偶次系数由 $a_0$ 定，奇次系数由 $a_1$ 定。取 $a_1=0$，就把两支解分开。
\end{frame}

\begin{frame}{$s=0$：余弦}
  \[
    a_{j+2}=-\frac{\omega^2}{(j+2)(j+1)}a_j.
  \]
  \[
    a_2=-\frac{\omega^2}{2!}a_0,\qquad
    a_4=\frac{\omega^4}{4!}a_0,\qquad
    a_6=-\frac{\omega^6}{6!}a_0.
  \]
  于是
  \[
    y=a_0\sum_{m=0}^{\infty}(-1)^m\frac{(\omega x)^{2m}}{(2m)!}
    =a_0\cos(\omega x).
  \]
\end{frame}

\begin{frame}{$s=1$：正弦}
  \[
    a_{j+2}=-\frac{\omega^2}{(j+3)(j+2)}a_j,\qquad a_1=0.
  \]
  \[
    y=a_0 x\left[
      1-\frac{(\omega x)^2}{3!}
       +\frac{(\omega x)^4}{5!}
       -\cdots
    \right]
    =\frac{a_0}{\omega}\sin(\omega x).
  \]
  两个根给出两个线性无关的解。通解是 $A\cos(\omega x)+B\sin(\omega x)$。
\end{frame}

\begin{frame}{展开中心可以搬家}
  常点若在 $x_0$，把幂次换成 $(x-x_0)$：
  \begin{equation*}
    y(x)=\sum_{j=0}^{\infty}a_j(x-x_0)^{s+j},\qquad a_0\neq 0.
  \end{equation*}
  判别方程和递推的形状不变，只是每一步的导数多一个因子。谐振子的系数是常数，换原点仍然得到正弦和余弦，只是宗量改成 $\omega(x-x_0)$。
\end{frame}

% =============================================================
\section{正则奇点：贝塞尔方程}

\begin{frame}{代入贝塞尔方程}
  \[
    x^2 y''+x y'+(x^2-n^2)y=0,
    \qquad
    y=\sum_{j=0}^{\infty}a_j x^{s+j}.
  \]
  $x^2 y''$、$x y'$ 和 $-n^2 y$ 都产生 $x^{s+j}$；$x^2 y$ 产生 $x^{s+j+2}$，比它们高两阶。
  \vspace{0.3em}
  最低次 $x^s$ 只含 $a_0$：
  \[
    a_0\bigl[s(s-1)+s-n^2\bigr]=0.
  \]
  判别方程是 $s^2-n^2=0$，根为 $s=\pm n$。下面先做 $s=n$，并设 $n\geq 0$。
\end{frame}

\begin{frame}{奇次系数为零}
  $x^{s+1}$ 的系数是
  \[
    a_1(s+1-n)(s+1+n)=0.
  \]
  $s=n$ 时两个因子是 $1$ 和 $2n+1$，都不为零，所以 $a_1=0$。递推每步跳两阶，所有奇次系数都是零。
\end{frame}

\begin{frame}{递推}
  把 $x^{s+j}$ 的系数收成
  \[
    a_j\,j(2n+j)+a_{j-2}=0.
  \]
  改记指标，
  \begin{equation*}
    \boxed{\displaystyle
      a_{j+2}=-\frac{a_j}{(j+2)(2n+j+2)},
      \qquad j\geq 0.}
  \end{equation*}
  偶数项可以一直写下去：
  \[
    a_{2m}=(-1)^m\frac{a_0\,n!}{2^{2m}\,m!\,(n+m)!}.
  \]
\end{frame}

\begin{frame}{第一类贝塞尔函数}
  \[
    y=a_0\sum_{m=0}^{\infty}(-1)^m
      \frac{n!}{2^{2m}\,m!\,(n+m)!}\,x^{n+2m}.
  \]
  取 $a_0=1/(2^n n!)$，这个级数定义为 $J_n(x)$：
  \begin{equation*}
    \boxed{\displaystyle
      J_n(x)
      =\sum_{m=0}^{\infty}(-1)^m
        \frac{1}{m!\,(n+m)!}
        \left(\frac{x}{2}\right)^{n+2m}.}
  \end{equation*}
  $n$ 不是整数时，$(n+m)!$ 理解成 $\Gamma(n+m+1)$。$s=-n$ 且 $n$ 不是整数，同样的步骤给出另一个解 $J_{-n}(x)$。
\end{frame}

\begin{frame}{整数阶时第二个根失败}
  $n$ 是非负整数时，$s=-n$ 的递推会走到分母为零的一步，级数无法按 $x^{-n}$ 开头继续下去。得到的解与 $J_n$ 成比例，不是第二个独立解。
  \vspace{0.4em}
  判别根之差为整数时，Frobenius 方法常常只给出一个解。第二个解要用降阶，一般会含 $\ln x$。
\end{frame}

% =============================================================
\section{第二个解}

\begin{frame}{朗斯基行列式}
  两个解 $y_1,y_2$ 的\textbf{朗斯基行列式}（Wronskian）是
  \begin{equation*}
    \Delta(x)
    =y_1 y_2'-y_2 y_1'
    =\begin{vmatrix}
      y_1 & y_2 \\ y_1' & y_2'
    \end{vmatrix}.
  \end{equation*}
  $\Delta\not\equiv 0$ 当且仅当两个解线性无关。
  \vspace{0.3em}
  它们都满足 $y''+py'+qy=0$，相减消去 $q$：
  \[
    y_1 y_2''-y_2 y_1''+p\,(y_1 y_2'-y_2 y_1')=0,
  \]
  也就是 $\Delta'+p\,\Delta=0$。
\end{frame}

\begin{frame}{阿贝尔公式}
  分离变量，$\mathrm{d}\Delta/\Delta=-p\,\mathrm{d}x$，
  \begin{equation*}
    \boxed{\displaystyle
      \Delta(x)=\Delta_0\exp\Bigl(-\int^{x}p(t)\,\mathrm{d}t\Bigr).}
  \end{equation*}
  知道一个解以后，
  \[
    \frac{\mathrm{d}}{\mathrm{d}x}\Bigl(\frac{y_2}{y_1}\Bigr)
    =\frac{\Delta}{y_1^2},
  \]
  所以
  \begin{equation*}
    \boxed{\displaystyle
      y_2(x)=y_1(x)\int^{x}\frac{\Delta(t)}{y_1(t)^2}\,\mathrm{d}t.}
  \end{equation*}
\end{frame}

\begin{frame}{例：从余弦回到正弦}
  谐振子 $p=0$，故 $\Delta$ 是常数。取 $y_1=\cos(\omega x)$，$\Delta_0=\omega$，
  \[
    y_2=\cos(\omega x)\int\frac{\omega}{\cos^2(\omega x)}\,\mathrm{d}x
    =\cos(\omega x)\cdot\tan(\omega x)
    =\sin(\omega x).
  \]
  和 Frobenius 的第二支一致。这里 $p=0$，积分里没有对数。
\end{frame}

\begin{frame}{根差为整数时的对数}
  贝塞尔方程 $p=1/x$，$\Delta\propto 1/x$。第一个解 $y_1\sim x^n$ 时，
  \[
    y_2=y_1\int\frac{\mathrm{d}x}{x\,y_1^2}.
  \]
  $n$ 是整数，被积函数含 $x^{-1}$ 或更低的整数幂，积分产生 $\ln x$，再乘上 $y_1$ 的级数。
  \vspace{0.35em}
  一般地，判别根之差为整数时，第二个解写成
  \[
    y_2=A\,y_1\ln|x-x_0|
    +\sum_{k=0}^{\infty}b_k(x-x_0)^{s_2+k}.
  \]
  $A=0$ 时对数项不出现，两个解都是纯级数。
\end{frame}

\begin{frame}{这一讲留下的公式}
  正则奇点的判据：$(x-x_0)p$ 与 $(x-x_0)^2 q$ 有限。
  \vspace{0.35em}
  Frobenius 级数 $\sum a_j(x-x_0)^{s+j}$，最低次决定 $s$。谐振子给出 $\cos(\omega x)$ 与 $\sin(\omega x)$；贝塞尔方程给出
  \[
    J_n(x)
    =\sum_{m=0}^{\infty}(-1)^m
      \frac{1}{m!\,(n+m)!}
      \left(\frac{x}{2}\right)^{n+2m}.
  \]
  只有一个级数解时，
  \[
    y_2=y_1\int\frac{\Delta}{y_1^2}\,\mathrm{d}x,
    \qquad
    \Delta=\Delta_0 e^{-\int p\,\mathrm{d}x}.
  \]
\end{frame}

\end{document}
